% Transcription — fonds Alexandre Grothendieck, shelfmark 7, batch 1 (pages 1-20).
% Established from the facsimile archives/batches/7/batch-01.pdf, cut from a
% copy of 7.pdf fetched through the project's own relay
% (https://grothendieck-relay.onrender.com/source/7.pdf) rather than from
% grothendieck.umontpellier.fr directly; per the skill transcribe-grothendieck.
% The batch file carries no cover sheet: PDF page k is the archivists' page k,
% and their pencil numbers bottom-left agree throughout. The folder is
% seventy-one pages; this is the first of four batches (the others are
% transcribed in parallel by other passes).
%
% Pass: Opus 5.5 (claude-opus-5-5), 2026-09-26 — first pass, unchecked against the pages by a human.
%
% The folder sits in the inventory group "4-9" « Cristaux »; it has its own
% date in the catalogue, which \dating{} carries verbatim, with no group
% sentence.
%
% Pages skipped: 1 and 19, blank cream sheets carrying nothing but
% show-through from the neighbouring leaves (no word of his on either).
%
% Pages summarised: none. Everything from 2 to 18 and page 20 is his own
% manuscript and is transcribed in full, as far as the hand allows.
%
% His pagination and dates. None of the leaves carries a page number of his.
% Two runs, each dated by him at the top right in pencil:
%   2-18  « Juillet-Août 1967 » (the third digit is barely formed; read 1967,
%         flagged \uncertain), titled « Anneaux simpliciaux canoniques et
%         modules stratifiés », with numbered paragraphs 1-6 (pp. 2, 4, 6-7,
%         10, 13), then, on p. 16, an unnumbered « § -- Commentaires sur la
%         définition d'une cohomologie p-adique en car. p > 0 » running to
%         p. 18 (« Tests décisifs »). Blue ink from p. 6 on.
%   20    « Nov. Déc. 196? » — the last digit is written over in red and is
%         not read with certainty (the catalogue title says novembre-décembre
%         1967); a new run numbered « 1. » at the top, « Cohomologie de De
%         Rham des var. alg. », which continues into batch 2.
%
% What the batch is about. For X -> S a scheme, the formal completions
% X^{[n]^} of X^{[n]} = X x_S ... x_S X along the diagonal give a
% cosimplicial ring P^* (principal parts), P^0 = O_X; the shifted
% cosimplicial ring Q^* = SP^*, (SP)^i = P^{i+1}, with alpha : P^* -> (SP)^*
% (for X = Spec A over k: Gamma P^i = A^{(x)(i+1)}, alpha^i(x_0...x_i) =
% x_0 (x) ... (x) x_i (x) 1). For a Module M, P^i[M] = P^i (x) M via the
% extreme-right structure (pro-objects in general), Q^*[M] a cosimplicial
% Module over Q^*; stratifications as descent data; Q^0[M] is the
% pro-Module of principal parts of M, stratified, functor in Mod(X/S)
% (Artin-Rees). Then H^*(X,M) = H^*(X, lQ^*[M]) with a spectral sequence,
% a limit formula, the extension to complexes M^. in D^+(X/S), the
% « formalisation » of complexes of differential operators; §5 the
% sheaves V^q = H^q(lQ^0[M^.]) as stratified Modules and the spectral
% sequence H^p_str(X, V^q) => H^*(X, M^.), with the corollary for V^q = 0,
% q != q_0; §6 the De Rham complex Omega^.(V) of a stratified Module and
% H^*_DR(X/S, V) = H^*_str(X/S, V) for X/S smooth (« lemme de Poincaré
% formel »). Pp. 16-18: comments on crystalline cohomology over W(k) and
% its torsion defects, and his proposal to enlarge the crystalline site by
% triples (U_0, W, i : U_0 -> W) with W « affine formal » and a nilpotent
% ideal, with « tests décisifs » a) and b). Page 20 opens a talk on the De
% Rham cohomology of algebraic varieties: 1) differentiable manifolds,
% 2) complex analytic manifolds (Hodge-to-De Rham spectral sequence),
% 3) algebraic varieties and the comparison with H^*(X^an, C).
%
% Notation fixed once for the batch: \mathcal{P} for his script P (sheaves
% of principal parts), \mathcal{Q} for the curly Q of pp. 5-14 (= SP);
% X^{[n]} and X^{\widehat{[n]}} as he writes the brackets; \mathbb{H} for
% his double-stroked H (hypercohomology), \mathcal{H} for his script H
% (cohomology sheaves); \underline{V} underlined as written (a stratified
% Module); \mathcal{O}_X for his O_X, whether or not he underlines it.
%
% Hardest pages and readings to check first: the connective prose of pp.
% 2-4 and 6-8 (fast drafting in black ink; formulas carry, sentences often
% do not); the clause after « exceptionnelles » on p. 4; « (Artin-Rees) » on
% p. 6; the circled margin of p. 9 (« Mention ... +1 ... »); the dates on
% pp. 2 and 20.

\documentclass[11pt,a4paper]{article}
% Le préambule commun à toutes les transcriptions du fonds Grothendieck.
%
% Il tient en une idée : **l'incertitude se compose, elle ne se résout pas.**
% Une transcription qui lisse un mot douteux a détruit la seule chose pour
% laquelle on la faisait — un lecteur doit pouvoir savoir, à chaque endroit,
% ce qui était sur le feuillet et ce qui a été ajouté. Les six macros qui
% suivent sont donc l'appareil critique, et non de la décoration.
%
% Les mêmes commandes sont reconnues par `scripts/render.mjs`, qui produit la
% vue de lecture du site. Une macro ajoutée ici doit l'être aussi là-bas,
% faute de quoi le rendu échouera bruyamment — ce qui est voulu : un rendu
% silencieusement faux est pire qu'un rendu absent.

\ProvidesPackage{grothendieck}[2026/08/08 Transcriptions du fonds Grothendieck]

\usepackage{amsmath,amssymb,amsthm}
\usepackage{mathrsfs}
\usepackage{stmaryrd}
% Les diagrammes commutatifs sont partout dans ce fonds : c'est la langue
% même de la géométrie algébrique de Grothendieck, et une transcription qui
% les rendrait en prose ne transcrirait rien.
\usepackage{tikz-cd}
% Français par défaut — c'est la langue des feuillets — mais l'anglais est
% chargé aussi : la traduction et le résumé sont des documents anglais, et la
% typographie française (espaces avant les deux-points) y serait une faute.
% Ces documents basculent par \selectlanguage{english} en tête de corps.
\usepackage[english,french]{babel}
\usepackage{fontspec}
\usepackage{geometry}
\geometry{a4paper,margin=25mm,marginparwidth=28mm}
\usepackage{marginnote}
\usepackage{xcolor}
% ulem, not soul. `soul`'s \st and \ul are built on a character-by-character
% scan that XeLaTeX's font machinery breaks: the document fails with an
% undefined control sequence rather than degrading. ulem does the same two jobs
% and is Unicode-engine safe. `normalem` stops it hijacking \emph, which the
% transcriptions use for Grothendieck's own emphasis.
\usepackage[normalem]{ulem}
\usepackage{hyperref}
\hypersetup{colorlinks=true,linkcolor=black,urlcolor=black,pdfborder={0 0 0}}

% --- Métadonnées du lot -----------------------------------------------------
% Reprises telles quelles par le script de rendu, qui les affiche en tête de
% la vue de lecture. Elles fixent ce que le fichier prétend couvrir : sans
% elles, une transcription flotte sans point d'attache dans un fonds de seize
% mille feuillets.

\newcommand{\folder}[1]{\gdef\grothFolder{#1}}
\newcommand{\batch}[1]{\gdef\grothBatch{#1}}
\newcommand{\leaves}[2]{\gdef\grothFirst{#1}\gdef\grothLast{#2}}
\newcommand{\foldertitle}[1]{\gdef\grothFoldertitle{#1}}
\gdef\grothFolder{?}\gdef\grothBatch{?}\gdef\grothFirst{?}\gdef\grothLast{?}\gdef\grothFoldertitle{}

% --- Le résumé --------------------------------------------------------------
%
% Il ouvre la lecture modernisée, sous le titre « Résumé », et il est composé
% en petit. Ce n'est pas un ornement : c'est le seul passage du document qui
% ne soit pas des mathématiques, et un lecteur doit pouvoir le reconnaître
% avant de l'avoir lu — puis le sauter s'il connaît déjà le sujet, ou s'y
% arrêter s'il ne le connaît pas. Le filet qui le suit marque la reprise du
% corps.

\newenvironment{resume}
  {\small\setlength{\parskip}{0.45em}}
  {\par\medskip\hrule\medskip}

% --- Mots-clés ---------------------------------------------------------------
%
% En anglais, en clôture du résumé de la lecture modernisée : le vocabulaire
% moderne sous lequel chercher ce que les feuillets construisent. Le manifeste
% du site (`npm run manifest`) les extrait pour étiqueter la cote entière —
% cette ligne est la source unique des tags, il n'y a pas de fichier à côté.
\newcommand{\keywords}[1]{\par\smallskip\noindent
  {\footnotesize\textsc{Keywords} --- \textit{#1}}\par}

% --- L'appareil critique ----------------------------------------------------

% Le numéro de feuillet, en marge. C'est l'ancre qui permet de régler un
% désaccord sur une formule en regardant *un* feuillet plutôt qu'une chemise
% de six cents. Le site s'en sert aussi pour tourner les pages du fac-similé
% quand on fait défiler la transcription.
\newcommand{\leaf}[1]{%
  \marginnote{\footnotesize\color{gray!70}\textsf{#1}}%
  \label{leaf:#1}\ignorespaces}

% Illisible : on ne devine pas. Un mot inventé qui se lit comme les autres est
% le pire résultat possible.
\newcommand{\ill}{\textcolor{red!60!black}{[\,\ldots\,]}}

% Lecture douteuse : le mot est donné, et signalé comme incertain.
\newcommand{\uncertain}[1]{\uline{#1}}

% Ajout de l'éditeur — une lettre manquante, un mot sous-entendu.
\newcommand{\add}[1]{\textcolor{blue!50!black}{[#1]}}

% Biffé par Grothendieck lui-même. Ce qu'il a rayé fait partie du document :
% une rature dit souvent où la pensée a changé de direction.
\newcommand{\struck}[1]{\sout{#1}}

% Note du transcripteur, sur l'état matériel du feuillet.
\newcommand{\note}[1]{\par\smallskip{\small\color{gray!60!black}\textit{[#1]}}\par\smallskip}

% Note marginale de Grothendieck, distincte du corps du texte.
\newcommand{\marginal}[1]{\par\smallskip\noindent\hspace{1em}%
  {\small\color{gray!60!black}#1}\par\smallskip}

% --- Titre ------------------------------------------------------------------

\AtBeginDocument{%
  \begin{center}
    {\large\bfseries Fonds Alexandre Grothendieck --- cote n\textsuperscript{o}~\grothFolder}\\[2pt]
    {\itshape\grothFoldertitle}\\[4pt]
    {\small feuillets \grothFirst--\grothLast\ (lot \grothBatch)}\\[2pt]
    {\footnotesize\color{gray!60!black}Transcription établie d'après le fac-similé numérisé
     par l'Université de Montpellier.\\
     Les lectures incertaines sont soulignées, les illisibles notées [\,\ldots\,],
     les ajouts entre crochets.}
  \end{center}
  \vspace{1em}}

\endinput


\folder{7}
\batch{1}
\pages{1}{20}
\dating{1966-1967}
\watermark{Édition de démonstration}
\foldertitle{Notes laïus novembre-décembre 1967 (Cristaux) : notes manuscrites (1966-1967)}

\begin{document}

\page{2}

\marginal{Juillet-Août \uncertain{1967}}\note{au crayon, en haut à droite ;
le troisième chiffre de l'année est à peine formé}

\section*{Anneaux \struck{\ill{}} simpliciaux canoniques et modules stratifiés}

\struck{\ill{}} [Formellement, les développements qui suivent sont
indépendants des \struck{\ill{}} histoires de sites cristallins ou
\uncertain{stratifiants}].

1. Soit $X \to S$ un préschéma relatif. Pour tout entier $n \geq 0$, soit
\[
  X^{[n]} = \underbrace{X \times_S \cdots \times_S X}_{n+1},
\]
considéré comme objet \struck{\ill{}} de $(\mathrm{Sch})_{/S}$ \ill{} \ill{}
\struck{diag} l'immersion \struck{diag} $X \to X^{[n]}$\note{le mot « diag »,
d'abord écrit, est biffé deux fois, dans la phrase et au-dessus de la
flèche}. Soit $X^{\widehat{[n]}}$ le complété formel \ill{} le long de $X$.
Pour $n \geq 0$ variable, on \ill{} un objet simplicial : \ill{} des
\ill{} un objet simplicial \ill{} aux $X^{\widehat{[n]}}$, \ill{} dans
$\mathrm{Sch}f/X$, les espaces sous-jacents étant d'ailleurs tous égaux à
$\mathrm{esp}(X)$. Soit $\mathcal{P}^n$ \struck{\ill{}} le faisceau
structural de $X^{\widehat{[n]}}$, \uncertain{Ann}-pro-objet de \struck{\ill{}}
\struck{\ill{}} considéré comme \ill{} de la catégorie des faisceaux d'anneaux
$\mathrm{Pro}(\mathrm{Ann}(\mathrm{esp}(X)))$ \ill{} $\mathrm{esp}(X)$,
\ill{} Pour $n$ variable, on trouve ainsi un Anneau cosimplicial
$\mathcal{P}^{*}$ sur $\mathrm{esp}(X)$, avec
\[
  \mathcal{P}^{0} = \mathcal{O}_X .
\]
Faire mention que \emph{\uncertain{certains}}

\page{3}

\ill{} un $\mathcal{O}_X$-Algèbre cosimpliciale, ainsi les deux \ill{}
$\mathcal{O}_X = \mathcal{P}^{[0]} \rightrightarrows \mathcal{P}^{[1]}$ sont en
général distincts, et mettent sur $\mathcal{P}^{[1]}$ deux structures
\emph{distinctes} de $\mathcal{O}_X$-algèbres.

Considérons l'objet \uncertain{décalé} $\mathcal{SP}$, défini par
\[
  (\mathcal{SP})^{i} = \mathcal{P}^{i+1}
\]
\ill{} fonctoriellement \ill{}
\[
  (\mathcal{SP})(I) = \mathcal{P}(I \sqcup e)
\]
\ill{} un ensemble \ill{} à un élément. L'homomorphisme \add{fonctoriel}
$I \to I \sqcup e$ définit $\mathcal{P}(I) \to \mathcal{P}(I \sqcup e)$,
i.e.
\[
  \alpha^{*} : \mathcal{P}^{*} \to (\mathcal{SP})^{*}
  \qquad \text{i.e.}\quad \alpha^{i} : \mathcal{P}^{i} \to \mathcal{P}^{i+1} .
\]
Si p.\,ex.\ $S = \operatorname{Spec} k$ ($k$ un anneau), $X = \operatorname{Spec} A$,
$X$ étant tel que $X \to X \times_S X$ soit une immersion nil-\ill{}
[p.\,ex.\ $X$ radiciel et \ill{} sur $S$], alors, \ill{} aux sections,
\[
  \Gamma \mathcal{P}^{i} = \bigotimes_{k}^{i+1} A
  = \underbrace{A \otimes_k A \otimes \cdots \otimes_k A}_{i+1},
\]
les opérations simpliciales \ill{} \ill{} \ill{}, et l'homomorphisme $\alpha$
étant donné par
\[
  \underbrace{A \otimes \cdots \otimes A}_{i+1}
  \longrightarrow
  \underbrace{A \otimes \cdots \otimes A}_{i+2},
  \qquad
  \alpha^{i}(x_0 \otimes \cdots \otimes x_i) = x_0 \otimes \cdots \otimes x_i \otimes 1 .
\]

\page{4}

2. Soit \struck{$(M, \ill{})$} un Module sur $X$. \struck{stratifié} \ill{}
On pose pour $i \geq 0$
\[
  \mathcal{P}^{i}[M] = \mathcal{P}^{i} \otimes_{\mathcal{O}_X} M ,
\]
où $\mathcal{P}^{i}$ est considéré comme Algèbre sur $\mathcal{O}_X$ par sa
structure \emph{extrême-droite}. \struck{Donc $\mathcal{P}^{i}(M)$ est un}
Par définition, le produit tensoriel doit être interprété comme un
pro-objet [de même que $\mathcal{P}^{i}$] --- sauf dans des conditions
exceptionnelles \ill{} que \ill{} \ill{} $X$ loc.\ noeth.,
\ill{} [\ill{} $M$ \uncertain{lisse} de type fini, si $X$ \ill{}
$\Omega_{X/S}$ et $M$ de type fini sur $X$] qui \ill{} \ill{} \ill{}
d'interpréter les $\mathcal{P}^{i}$ et $\mathcal{P}^{i}[M]$ comme des
faisceaux au lieu de profaisceaux [car dans ces cas au lieu de
\ill{}, on remplace les profaisceaux \ill{} par leur \struck{\ill{}}
$\varprojlim$]. On notera que les $\mathcal{P}^{i}(M)$ est un
$\mathcal{P}^{i}$-Module, dépendant fonctoriellement de $M$ variable dans
$\mathrm{Mod}(\mathcal{O}_X)$. \ill{} on prend comme \ill{} de $M$ \ill{}

\page{5}

\ill{} $\mathcal{O}_X$ sur les $\mathcal{Q}^{i} = \mathcal{P}^{i+1}$
\ill{} \ill{} à cette structure. Il s'ensuit que les
\[
  \mathcal{Q}^{i}[M] = \mathcal{Q}^{i} \otimes_{\mathcal{O}_X} M
\]
[str.\ ext.\ dr.]\note{écrit sous le signe $\otimes$, avec un trait qui l'y
rattache : le produit tensoriel est pris pour la structure extrême-droite}
sont les composantes d'\ill{} \struck{\ill{}} cosimplicial, noté
$\mathcal{Q}^{*}[M]$, \ill{} \ill{} qui est en fait un \emph{Module}
cosimplicial sur l'anneau cosimplicial $\mathcal{Q}^{*}$. Ce
$\mathcal{Q}^{*}[M]$ dépend également de façon fonctorielle de $M$, quand
$M$ parcourt $\mathrm{Mod}(X/S)$.

3. Notons une conséquence \uncertain{importante} \ill{} de \ill{} des Modules
cosimpliciaux \ill{} $\mathcal{Q}^{*}[M]$ sur $\mathcal{P}^{*}$ : \struck{\ill{}}
\ill{} \ill{} $\mathcal{Q}^{i}[M]$ \ill{} \struck{\ill{}} c'est que les
\ill{} \ill{} \ill{} autres \ill{} \ill{} \ill{} \ill{} les opérations
\ill{} \ill{} changement de base \ill{} les \uncertain{opérations}
simpliciales $\mathcal{P}^{i} \to \mathcal{P}^{j}$ [\ill{} \ill{} \ill{}
\struck{\ill{}} \ill{} \ill{} d'un $\varprojlim$, non seulement

\page{6}

d'un [~~]. Formellement, [de façon précise \ill{} cela plus \ill{} \ill{}
des pro-Modules \ill{} \ill{} de Modules \ill{} \ill{}] on voit que
$\mathcal{Q}^{0}[M]$ est un (pro)Module stratifié sur $X$ rel $S$, et les
$\mathcal{Q}^{i}[M]$ sont les composantes des (pro)Modules \struck{\ill{}}
cosimplicial correspondant.\marginal{\ill{} \ill{} \ill{}}\note{le
passage « on voit que \ldots{} correspondant » est marqué d'un double trait
vertical dans la marge gauche ; l'ajout au-dessus de la ligne suivante,
entouré et relié à $\mathcal{Q}^{0}[M]$, n'est pas lu}
Ce pro-Module stratifié $\mathcal{Q}^{0}[M]$ \ill{} \ill{} fonctoriellement
de $M$ parcourant $\mathrm{Mod}(X/S)$ [Modules sur $X$ \ill{}
\uncertain{rel $S$} \ill{} \emph{\uncertain{ordre} borné}]\note{« rel $S$ »
est écrit au-dessus d'un mot biffé ; « ordre borné » est souligné}, on
\ill{} \ill{} \ill{} \ill{} \ill{} comme \ill{} \ill{} \ill{}
\uncertain{Artin-Rees}. De plus, on \uncertain{récupère} $M$ à partir de
$\mathcal{Q}^{0}[M]$ [\ill{} : stratification] \ill{} \ill{} \ill{} \ill{}
le sous-faisceau des \emph{sections \ill{}} [\ill{} la
structure\note{le passage « Ce pro-Module stratifié \ldots{} la structure »
est lui aussi marqué d'un trait vertical dans la marge gauche}

\page{7}

de $\mathcal{O}_X$-Module \ill{} de l'extrême droite [N.B. Les structures
extrême-droites des $\mathcal{Q}^{i}[M]$ sont compatibles avec les op.\
simpliciales, i.e.\ $\mathcal{Q}^{i}[M]$ est un Module cosimplicial sur
$\mathcal{O}_X$ \ill{} \ill{}, \ill{} \ill{} \ill{} extrême-droites ; mais
\ill{} \ill{} \ill{} \ill{} \ill{} structures \ill{} \ill{} \ill{}
conservée par les \ill{} $M \to N$ qui \ill{} \ill{} $\mathcal{O}_X$-\ill{}
(\ill{} \ill{} \ill{} \ill{} diff.\ rel $S$ \ill{})]. \ill{} \ill{} un
\ill{} $\mathrm{Mod}(X/S) \to$ \ill{} $D : M \to N$ \ill{}
$\mathcal{Q}^{0}[D] : \mathcal{Q}^{0}[M] \to \mathcal{Q}^{0}[N]$ \ill{}
\ill{} \ill{} \ill{} comme l'\ill{} \ill{} \ill{} les
\uncertain{sous}-faisceaux des \struck{\ill{}} \add{sections}
\uncertain{horizontales}.

4. Une façon \uncertain{équivalente} d'obtenir $\mathcal{Q}^{*}[M]$ \ill{}
de considérer
\[
  X^{\widehat{[2]}} \rightrightarrows X^{\widehat{[1]}}
  \xrightarrow{\widehat{\mathrm{pr}}_2} X^{\widehat{[0]}} = X
\]
et les \uncertain{images} \ill{}, [dans les cat.\ \ill{} \ill{} formels sur
$X$], \struck{\ill{}} c'est ainsi

\page{8}

\ill{} \ill{} l'anneau \add{cosimplicial} $\mathcal{Q}^{*}$, et
\[
  \mathcal{Q}^{*}[M] = \mathcal{Q}^{*} \otimes_{\mathcal{O}_X} M .
\]
Comme $X^{\widehat{[1]}}$ \ill{} \ill{} section sur $X^{\widehat{[0]}}$, on
obtient aussitôt que \emph{$\mathcal{Q}^{*}[M]$ est une résolution de $M$}
[\ill{} \ill{} \ill{} \ill{} les $\mathcal{Q}^{i}[M]$ les structures
extrêmes droites de $\mathcal{O}_X$-Modules]. Il faut \ill{} \ill{}
cependant \ill{} \ill{} \ill{} \ill{} \ill{} \ill{} \ill{} cette expression
[les A.R.-pro-modules \ill{} \ill{} catégorie abélienne\,?\,\ldots]
\ill{}-il \ill{} \ill{} \ill{} faut \ill{} \ill{} \ill{} \ill{} \ill{}
hypothèse de quasi-cohérence sur $M$. En tout cas, \struck{\ill{}}
\uncertain{désignant} par \ill{} \ill{} \ill{} le \ill{} \ill{}
$\varprojlim$ des \ill{}, \ill{} \ill{} projets, on trouve que
$l\mathcal{Q}^{*}[M]$ est une résolution de $M$ \ill{} $M$ \ill{}.

\page{9}

et par suite on trouve un bon cas
\[
  H^{*}(X, M) \simeq \mathbb{H}^{*}(X, l\mathcal{Q}^{*}[M])
  \Longleftarrow
  E_2^{pq} = H^{p}\bigl(i \mapsto H^{q}(X, l\mathcal{Q}^{i}[M])\bigr)
\]
[Sous des conditions convenables, \ill{} aura
\[
  H^{q}(X, l\mathcal{Q}^{i}[M]) = \varprojlim_{n}
  H^{q}(X, \mathcal{Q}^{i}_{\uncertain{n}}[M]_{(n)})
  \;\bigl(\simeq H^{q}(X^{\widehat{[i+1]}}_{n}, M)\bigr)
\]
\ill{} \ill{} \ill{}
\[
  \mathcal{Q}^{i}[M] = \varprojlim_{n} \mathcal{Q}^{i}[M]_{(n)}, \quad
  \mathcal{Q}^{i}[M]_{(n)} = \mathcal{Q}^{i}_{(n)} \otimes_{\mathcal{O}_X} M, \quad
  \mathcal{Q}^{i}_{(n)} = \mathcal{P}^{i+1}_{(n)} = \cdots\;]
\]
\note{la parenthèse $(\simeq H^{q}(X^{\widehat{[i+1]}}_n, M))$ est ajoutée
au-dessus de la ligne et entourée ; en bout de ligne, un
« $= H^{q}(X^{[i]}, M)$ » est biffé}\marginal{Mention \ill{} $+1$ \ill{}\,!
\struck{$X^{[i+1]}$} $= [X/S]^{i+2}$ \ill{} \ill{}\,!)}\note{dans la marge
gauche, entouré ; le « $+1$ » est lui-même entouré}

Si \ill{} complexe $M^{\cdot}$ \ill{} \ill{} \emph{\ill{}} $M^{\cdot}$ de
\struck{\ill{}} $D^{+}(\mathrm{Mod}(X/S)) = D^{+}(X/S)$, on obtient un
bicomplexe $\mathcal{Q}^{*}[M^{\cdot}]$, tel que $l\mathcal{Q}^{*}[M^{\cdot}]$
est une résolution de $M^{\cdot}$, et par suite
\[
  \mathbb{H}^{*}(X, M^{\cdot}) \simeq \mathbb{H}^{*}(X, \mathcal{Q}^{*}[M^{\cdot}])
  \simeq \mathbb{H}^{*}(X_{\mathrm{str}}, \mathcal{Q}^{0}[M^{\cdot}]) .
\]
[$\mathcal{Q}^{0}[M^{\cdot}]$ s'\uncertain{appelle} la
\emph{formalisation} des complexes \struck{\ill{}} d'op.\ différentiels
$M$.]\note{ce crochet est marqué d'un triple trait vertical dans la marge}

\page{10}

5. Soit \uncertain{toujours} $M^{\cdot} \in \mathrm{ob}\, D^{+}(X/S)$, et
considérons le complexe $\mathcal{Q}^{0}[M^{\cdot}]$ \struck{de}
\struck{$l\mathcal{Q}^{0}[$} de (pro-\struck{$\mathcal{O}_X$-Modules})
à stratifications, \ill{} \ill{} \ill{} \ill{} \ill{} \ill{} que celle de
$M^{\cdot}$ : \ill{} \ill{} \ill{} \ill{} \ill{} complexe \ill{} additif
d'abord, en \ill{} \ill{} \ill{} faisceaux de sections
\uncertain{Artiniennes}, puis \ill{} \ill{} $\mathcal{O}_X$-structures sur
les $M^{\cdot}$ comme \ill{} des structures extrêmes droites des
$\mathcal{Q}^{0}[M^{i}]$ (lesquelles, rappelons-le, \ill{}, ne sont
cependant pas respectées par l'op.\ diff.\ de $\mathcal{Q}^{0}[M^{\cdot}]$).
Considérons les $\mathcal{H}^{i}(l\mathcal{Q}^{0}[M^{\cdot}])$, et supposons
que leur formation commute aux changements de base \struck{\ill{}}
cosimpliciaux $\mathcal{P}^{0} \to \mathcal{P}^{\uncertain{n}}$ ; ceci
\ill{} \ill{} \ill{} par des \uncertain{hypothèses}

\page{11}

de \add{type} platitude sur les $\mathcal{Q}^{0}[M^{\cdot}]$ et les
$\mathcal{H}^{i}(\mathcal{Q}^{0}[M^{\cdot}])$, \struck{\ill{}} \ill{} (plus
économiques, \ill{} \ill{}, \ill{} que \ill{} ne \ill{} \ill{} \ill{}
\ill{} de platitude sur $\mathcal{O}_X$ de $\mathcal{Q}^{0}[M^{\cdot}]$
\ill{} \ill{} de $X^{\widehat{[1]}} \to X^{0}$] \ill{} un hypothèse
\ill{} « platitude \uncertain{infinitésimale} de $X/S$ », \ill{} \ill{}
\add{\uncertain{sur} platitude \ill{}}
$X^{\widehat{[1]}} \to X$ --- p.\,ex.\ $X/S$ \uncertain{lisse}
\ldots{} [critères : \uncertain{diagonale}]. Alors\note{ce passage est
marqué d'un double trait vertical dans la marge}
\begin{enumerate}
\item[1°)] Les $V^{i} = \mathcal{H}^{i}(l\mathcal{Q}^{0}[M])$ sont de
\ill{} \struck{\ill{}} façon naturelle des $\mathcal{O}_X$-Modules (par la
\supplied{$\mathcal{O}_X$-}structure \uncertain{produit}) \emph{stratifiés}.
\item[2°)] Dans le bicomplexe $l\mathcal{Q}^{*}[M^{\cdot}]$, prenant la
\add{\uncertain{première}} \struck{cohomologie} suivant l'op.\ diff.\
provenant de $M^{\cdot}$, \ill{} \ill{} \struck{\ill{}} complexes
\ill{} \ill{} \uncertain{fournissant} les stratifications
\struck{\ill{}} des $V^{i}$, $C^{\cdot}(V^{i})$ \ill{} \ill{}
\item[3°)] Par suite, on trouve une suite spectrale
\end{enumerate}

\page{12}

\[
  \mathbb{H}^{*}(X, M^{\cdot}) \simeq \mathbb{H}^{*}(X, \mathcal{Q}^{*}[M^{\cdot}])
  \Longleftarrow E_2^{pq} = \mathbb{H}^{p}(X, C^{\cdot}(V^{q}))
\]
(vérif.\ à \uncertain{EGA} $0_{\mathrm{III}}$\,?). On \ill{} \ill{}
l'écrire, en \ill{} la définition, \ill{} \ill{} Module stratifié
$\underline{V}$ relativement à $S$,
\[
  \mathbb{H}^{p}_{\mathrm{str}}(X/S, \underline{V})
  = \mathbb{H}^{p}(X, C^{\cdot}_{\mathrm{str}}(\underline{V})) ,
\]
\note{l'indice « str » est entouré et relié par un trait à la place d'un
« $X/S$ » biffé dans le premier membre}
sous la forme
\[
  \boxed{\;
  \mathbb{H}^{*}(X, M^{\cdot}) \Longleftarrow
  E_2^{pq} = \mathbb{H}^{p}_{\mathrm{str}}(X, V^{q}(M^{\cdot})),
  \quad\text{où}\quad
  V^{q}(M^{\cdot}) = \mathcal{H}^{q}(\mathcal{Q}^{0}[M^{\cdot}])
  \;}
\]

\emph{Corollaire}. Supposons que l'on ait
\[
  V^{q}(M^{\cdot}) = 0 \quad \text{pour } q \neq q_0 ,
\]
alors on \ill{} \ill{} des isom.\ canoniques
\[
  \mathbb{H}^{n}(X, M^{\cdot}) \simeq
  \mathbb{H}^{n-q_0}_{\mathrm{str}}(X, V^{q_0}(M^{\cdot})) .
\]
Ceci montre, (\add{supposant p.\,ex.\ que $q_0 = 0$,}\note{ajouté au-dessus
de la ligne}) que le premier \ill{}

\page{13}

ne dépend que du Module stratifié $\underline{V} = V^{q_0}(M^{\cdot})$ et de
\struck{\ill{}} $q_0$, mais pas du complexe différentiel particulier choisi
pour « résoudre » $\underline{V}$.

6. Or, \struck{\ill{}} \uncertain{Soit} \ill{} \struck{\ill{}} \ill{} \ill{}
\ill{} d'un Module stratifié $\underline{V}$ sur $X$ rel $/S$, \ill{}
\ill{} \ill{} lui associer de façon naturelle un \emph{complexe de De Rham},
\ill{} \ill{} l'objet gradué sous-jacent est
\[
  \Omega^{\cdot}_{X/S} \otimes \underline{V} = \Omega^{\cdot}(\underline{V})
\]
et on vérifie, \emph{lorsque} (\add{$X$ est lisse sur $S$ et} $S$ \ill{} de
\emph{\ill{}})\note{la parenthèse est ajoutée au-dessus de la ligne}, que ce
complexe \emph{\ill{}}, que \ill{} \ill{} « résolution différentielle » \ill{}
bien \ill{} de $\underline{V}$. Dans ce cas, on trouve donc une
interprétation \ill{} de la\marginal{« Lemme de Poincaré formel »}\note{dans
la marge gauche, entouré, avec un triple trait vertical ; plus bas, un
second mot entouré, « différentielle », à demi couvert d'une tache d'encre}

\page{14}

coh.\ de De Rham de $X/S$
\[
  H^{*}_{\mathrm{DR}}(X/S, \underline{V}) \simeq H^{*}_{\mathrm{str}}(X/S, \underline{V})
\]
au \ill{} \ill{} \ill{} que il \ill{} plus question de \struck{cohomologie}
\add{formes} différentielles\,! Si on ne fait plus d'hyp.\ de lissité
\ill{} de caractéristique, il \ill{} \ill{} \ill{} \ill{} \ill{} un
\add{\uncertain{homomorphisme} \ill{}}
\uncertain{canoniques}
\[
  \underline{V} \to \mathcal{H}^{0}(l\mathcal{Q}^{0}[\Omega^{\cdot}(\underline{V})])
\]
\note{l'insertion « un homomorphisme \ldots{} », au-dessus de la ligne, est
reliée par une accolade à cette formule}
\ill{} \ill{} \ill{} complexes de Modules stratifiés
\[
  \underline{V} \longrightarrow l\mathcal{Q}^{0}[\Omega^{\cdot}(\underline{V})]
\]
\ill{} \ill{} \ill{} au \uncertain{Cstr} des \ill{} \ill{} \ill{}, un
\emph{hom.\ can.}
\[
  \boxed{\;
  H^{*}_{\mathrm{str}}(X, \underline{V}) \longrightarrow
  H^{*}_{\mathrm{DR}}(X/S, \underline{V})
  \;}
\]
\note{l'indice « str » est entouré et relié à la place d'un mot biffé après
$X$}

\page{15}

qui est dans un \ill{} \ill{} $X/S$ est différentiellement lisse et $S$ de
car.\ nulle.

Je doute que ce soit \ill{} \ill{} \ill{} (p.\,ex.\ une \struck{\ill{}}
courbe elliptique)\note{la parenthèse est ajoutée au-dessus de la ligne ;
« doute » est entouré, et un long trait le relie à la marge} \struck{\ill{}}
\ill{} $X$ est \emph{lisse et projectif} [sur $S$ spectre d'un corps (alg.\
clos \ill{} \ill{}) de car.\ $p > 0$, et $\underline{V} = \mathcal{O}_X$
avec stratification canonique. Dans ce cas particulier,] il n'est pas
\ill{} clair à priori que les $H^{i}_{\mathrm{str}}$ soient de dim.\ finie,
ni qu'ils \ill{} nuls pour $i$ grand. Situation à éclaircir\,!

\emph{Remarque}. L'intérêt de relier la cohomologie de De Rham : la
\struck{\ill{}} $H_{\mathrm{str}}$ \ill{} que celle-ci, s'interprète comme
\ill{} \ill{} tour, \ill{} \ill{} « site stratifiant » \ill{} cohomologie
d'un \ill{} \ill{} variante \ill{} convenable, \ill{} \ill{} conduire à
la « coh.\ \uncertain{p-adique} en car.\ $p$ » \ill{}\marginal{C'est
\uncertain{sûrement} faux\,! (calculs \ill{} \ill{} \ill{} $X$ \ill{}
\uncertain{abélien}\,!)}\note{dans la marge gauche, dans un cadre ; la
première ligne au crayon, les suivantes à l'encre bleue, avec un triple
trait vertical}

\page{16}

\section*{§\,-- Commentaires sur la définition d'une cohomologie
\uncertain{p-adique} en car.\ $p > 0$}

$X_0$ schéma sur un corps parfait $k$ de car.\ $p > 0$,
$S = \operatorname{Spec} W(k)$. Alors
\[
  H^{*}(X_0/S_{\mathrm{cris}}, \mathcal{O}_{X_0\,\mathrm{cris}})
\]
est un $W$-module gradué. \struck{\ill{}} Mon espoir initial était qu'il
soit \struck{un} « le » bon candidat pour une cohomologie \uncertain{p-adique}.
Une première \ill{} \ill{} en \ill{}, c'est que lorsque $X_0$ \add{est
lisse sur $k$, \ill{} bien \ill{}} [se relève en un $X$ \struck{\ill{}}
\add{\ill{}} lisse sur $k$, alors d'après l'exposé précédent, on trouve
ainsi
\[
  H^{*}((X_0/S)_{\mathrm{cris}}, \mathcal{O}) = \varprojlim_{n}
  H^{*}((X_n/S_n)_{\mathrm{cris}}, \mathcal{O})
\]
et j'espérais qu'à la limite, on trouverait (pour certains modules
\uncertain{tensoriels}) la même chose que
$\varprojlim H^{*}_{\mathrm{DR}}(X_n/S_n) \simeq H^{*}_{\mathrm{DR}}(X/S)$.
(Il y \struck{\ill{}} a un \ill{} en tous cas\,!). C'est faux
malheureusement. Par exemple si $X_0/k$ est un schéma abélien
« ordinaire » i.e.\ \ill{} du point de vue formel est un tore, on trouve
que
\[
  H^{i}((X_0/S)_{\mathrm{cris}}, \mathcal{O}) = 0 \quad \text{si } i > \dim X_0\,!
\]
(Pourtant le $\varprojlim$ des coh.\ de De Rham est excellent\,!)
\struck{\ill{}} Quand on regarde dans le \uncertain{cas}

\page{17}

affine, on constate que les ennuis proviennent du fait que le passage à la
limite projective ne commute pas à la localisation par $p$
\[
  W \longrightarrow W[1/p] = K\,! .
\]

Je propose d'imaginer l'expédient suivant, consistant à remplacer le site
cristallin par un site plus gros, formé avec des « épaississements » pas
forcément infinitésimaux : objets les triples
$(U_0, \mathcal{W}, i : U_0 \to \mathcal{W})$, où $U_0$ un ouvert affine de
$X_0$, $\mathcal{W}$ un \struck{\ill{}} affine \add{(d'anneau $A_0$)}
formel complet \ill{}, i.e.\ \struck{\ill{}} correspondant à une
$W$-alg.\ $B$ sur $W$, $U_0 \to \mathcal{W}$ une \ill{} de $W$-alg.\
\struck{\ill{}} un morphisme \struck{\ill{}} induisant un \struck{\ill{}}
\emph{isomorphisme} \emph{nilpotent} $U_0 \to \mathcal{W} \times_S S_0$
(Cela correspond à un $W$-hom.\ surjectif $B \to A_0$ tel que
$B/pB \to A_0$ soit à noyau nilpotent). Comme topologie, une variante
convenable de la zariskienne (en \ill{} \ill{} : \ill{} d'ensembles
\ill{} \uncertain{par} les traces des ouverts sur $X_0$).

\page{18}

\emph{Tests décisifs} :
a) Si $X_0$ lisse \add{affine} sur $k$, on doit retrouver la définition de
\uncertain{W.M.}\note{« affine » est ajouté au-dessus de la ligne}
b) Si $X_0$ lisse et propre, et se relève en $X$ lisse et propre, on doit
retrouver (au moins modulo torsion) la cohomologie de De Rham
$H_{\mathrm{DR}}(X/S)$.

Il s'agit d'une première \struck{\ill{}}. \struck{\ill{}} Même si elle
réussit, il faudra beaucoup l'\uncertain{accomplir}, on \ill{} désire une
théorie cohomologique \ill{} \ill{} le projet, et \ill{} \ill{}
\uncertain{puisse} \ill{} \ill{} coup la \ill{} des \ill{}. Il y \ill{}
\ill{} \ill{} \ill{} \ill{} perspective\,!

\page{20}

\marginal{Nov.\ Déc.\ 196\uncertain{6}}\note{au crayon, en haut à droite ;
le dernier chiffre est repassé à l'encre rouge et le chiffre sous-jacent
n'est pas lu avec certitude}

\section*{1. Cohomologie de De Rham des var.\ alg.}

\note{« 1. » est écrit au-dessus du titre}
[But : faire sentir un complexe de \uncertain{pbs} gravitant autour de la
coh.\ de De Rham des V.A., et proposer un point de vue qui \struck{\ill{}}
me semble permettre de \uncertain{les} grouper et de les étudier. D'où un
(gros) programme de travail proposé.]\note{ce paragraphe, à l'encre bleue,
est ajouté entre le titre et le §\,1, et encadré}

§\,1. \emph{Yoga cohomologie de De Rham}

1) \emph{Variétés différentiables}. Énoncé du th.\ de De Rham. Son
explication en termes de théorie des faisceaux : \struck{\ill{}}
$\Omega^{\cdot}_{X}$ est une \emph{résolution} de $\mathbb{C}_X$ (lemme de
Poincaré), et les $\Omega^{i}_{X}$ sont fins, donc à cohomologie nulle en
dim $> 0$ (N.B. $X$ est supposé \uncertain{séparé}\ldots).

2) \emph{Variétés analytiques complexes}. Ça \ill{} \ill{} tel quel pour
les Stein : lemme de Poincaré (local) valable, mais il faut
$H^{j}(X, \Omega^{i}_{X}) = 0$ pour $j > 0$. En général, il faut prendre
l'hypercohomologie, on trouve
\[
  H^{*}(X, \mathbb{C}) \overset{\text{12}}{=}
  \mathbb{H}^{*}(X, \Omega^{\cdot}_{X})
  \Longleftarrow E_1^{pq} = H^{q}(X, \Omega^{p}_{X})
\]
\note{« 12 » est écrit au-dessus d'une ligne verticale qui relie
$H^{*}(X,\mathbb{C})$ à $\mathbb{H}^{*}(X,\Omega^{\cdot}_X)$ écrit
au-dessous ; lecture du chiffre douteuse}

3) \emph{Variétés algébriques} (non singulières). Travaillons avec la
top.\ de Zariski. Considérons
\[
  \mathbb{H}^{*}(X, \Omega^{\cdot}_{X}) \qquad \text{(hypercohomologie)}
\]
est-il isomorphe à $H^{*}(X^{\mathrm{an}}, \mathbb{C})$\,? On a en tous cas
\[
  \mathbb{H}^{*}(X, \Omega^{\cdot}_{X}) \longrightarrow
  \mathbb{H}^{*}(X^{\mathrm{an}}, \Omega^{\cdot}_{X^{\mathrm{an}}})
  \simeq H^{*}(X^{\mathrm{an}}, \mathbb{C})
\]

\end{document}
